\documentclass[12pt, oneside, a4paper]{book}
\usepackage[top=2.5cm, bottom=2.5cm, left=3cm, right=2.5cm, includehead]{geometry}
\usepackage[pdftex]{graphicx}
\usepackage{setspace}
\usepackage[latin1]{inputenc}
\usepackage[cmex10]{amsmath}
\usepackage{amsfonts}
\usepackage{amssymb}
\usepackage{hyperref}
\usepackage{url}
\usepackage{enumitem}
\usepackage{array}
\usepackage{ragged2e}
\usepackage{mathptmx}   % Times New Roman font
\usepackage{parskip}
\usepackage{tocloft}
\setlist[1]{itemsep=-17pt}
\usepackage[usenames,dvipsnames,pdftex]{color}
\usepackage{fancyhdr}
\setlength{\headheight}{16pt}
\setlength{\headsep}{12pt}
\setlength{\cftbeforetoctitleskip}{-18pt}
\setlength{\cftbeforeloftitleskip}{-18pt}

\fancypagestyle{phdthesis}{%
  \fancyhf{}
  \lhead{\leftmark}
  \rhead{}
  \rfoot{\thepage}
  \renewcommand{\headrulewidth}{0pt}
}

\fancypagestyle{plain}{%
  \fancyhf{}
  \renewcommand{\headrulewidth}{0pt}
  \rfoot{\thepage}
}

\newcommand{\HRule}{\rule{\linewidth}{0.5mm}}


\begin{document}

% ── Title Page ─────────────────────────────────────────────────────────────────
\newgeometry{top=1cm, bottom=1cm, left=1cm, right=1cm}
\begin{titlepage}
\begin{center}

\vspace*{1.3cm}

{\fontsize{18}{20}\selectfont \textbf{Frontend Architecture and System Design Program}}\\[0.7cm]

{\fontsize{12}{14}\selectfont by}\\[0.7cm]

{\fontsize{14}{16}\selectfont \textbf{Md. Miftahul Islam}}\\
{\fontsize{12}{14}\selectfont 221014030}\\[0.3cm]

{\fontsize{14}{16}\selectfont \textbf{Md. Nazirul Islam}}\\
{\fontsize{12}{14}\selectfont 221014029}\\[1.3cm]

\vspace{1cm}

{\fontsize{11}{13}\selectfont \textit{Internship report (CSE4099-A) submitted in partial fulfillment of the\\
requirements for the degree of}}\\[0.8cm]

{\fontsize{14}{16}\selectfont \textbf{Bachelor of Science in Computer Science and Engineering}}\\

\vspace{1cm}

{\fontsize{12}{14}\selectfont Under the supervision of}\\[0.7cm]

{\fontsize{13}{15}\selectfont \textbf{Dr. Muhammad Golam Kibria, SMIEEE}}\\
{\fontsize{11}{13}\selectfont Professor \& Head (Academic)}\\
{\fontsize{11}{13}\selectfont Department of Computer Science and Engineering}\\[0.8cm]

{\fontsize{13}{15}\selectfont \textbf{Md. Rafshanul Hoque Siam}}\\
{\fontsize{11}{13}\selectfont \textbf{Lead Software Engineer --- conneqtedagents.ai}}\\[1.8cm]

\includegraphics[width=0.20\textwidth]{logo.png}\\[1cm]

{\fontsize{12}{14}\selectfont \textbf{DEPARTMENT OF COMPUTER SCIENCE AND ENGINEERING}}\\[0.5cm]

\vspace{0.8cm}
{\fontsize{14}{16}\selectfont \textbf{UNIVERSITY OF LIBERAL ARTS BANGLADESH}}\\

\vfill
{\fontsize{10}{12}\selectfont \textbf{SPRING 2026}}

\end{center}
\end{titlepage}
\restoregeometry

% ── Copyright Page ─────────────────────────────────────────────────────────────
\newpage
\thispagestyle{empty}

\vspace*{\fill}
\begin{center}
\textcopyright{} Nazirul Islam \& Miftahul Islam \\[0.5cm]
All rights reserved
\end{center}
\vspace*{\fill}

% ── Declaration Page ───────────────────────────────────────────────────────────
\newpage
\thispagestyle{empty}

\begin{center}
{\large \textbf{DECLARATION}}

\vspace{2.0cm}

\begin{tabular}{>{\bfseries}p{0.34\textwidth} p{0.58\textwidth}}
  Project Title          & Frontend Architecture and System Design Program \\[4pt]
  Author                 & \textit{Md. Miftahul Islam}                     \\[2pt]
  Student ID             & 221014030                                       \\[4pt]
  Author                 & \textit{Md. Nazirul Islam}                      \\[2pt]
  Student ID             & 221014029                                       \\[4pt]
  Supervisor (Academy)   & Dr. Muhammad Golam Kibria, SMIEEE               \\[2pt]
  Supervisor (Industry)  & Md. Rafshanul Hoque Siam                        \\
\end{tabular}

\vspace{0.25cm}
\hrule
\vspace{1.2cm}

\setstretch{1.35}
\justifying
We declare that this internship report entitled
\textit{Frontend Architecture and System Design Program}
is the result of our own work except as cited in the references.
The internship report has not been accepted for any degree and is not
concurrently submitted in candidature of any other degree.
\setstretch{1.0}

\vspace{2.5cm}

\rule{0.60\textwidth}{0.4pt}

\vspace{0.25cm}

\begin{flushleft}
  \textbf{Md. Miftahul Islam \& Md. Nazirul Islam}\\[2pt]
  \textbf{221014030 \& 221014029}
\end{flushleft}

\vspace{0.7cm}

Department of Computer Science and Engineering\\
University of Liberal Arts Bangladesh

\vspace{0.5cm}

\textbf{Date:} April 22, 2026

\end{center}

% ── Letter of Transmission ─────────────────────────────────────────────────────
\newpage
\newgeometry{top=2.5cm, bottom=2.5cm, left=4cm, right=2.5cm}
\thispagestyle{empty}

\section*{Letter of Transmission}

\noindent April 22, 2026 \\[0.5cm]

\noindent The Chairperson,\\
Department of Computer Science \& Engineering,\\
University of Liberal Arts Bangladesh (ULAB),\\
Dhanmondi, Dhaka.\\[0.5cm]

\noindent Through: Supervisor, Dr.\ Muhammad Golam Kibria, SMIEEE\\[0.5cm]

\noindent \textbf{Subject: Submission of Internship Report on ``Frontend Architecture and System Design Program''}\\[0.5cm]

We are submitting herewith our internship report on \textit{Frontend Architecture and
System Design Program}. We had the opportunity to complete our internship at
\textbf{conneqtedagents.ai} under the industry supervision of \textbf{Md.\ Rafshanul
Hoque Siam}, Lead Software Engineer. Our internship spanned the Spring 2026 semester,
during which we actively participated in the design and development of scalable frontend
systems for a production web application.

It was a great pleasure for us to work at conneqtedagents.ai. Under the guidance of our
supervisors and in fulfilment of the requirements of the University of Liberal Arts
Bangladesh, this internship was carried out with full dedication and sincerity.

We have prepared this report with our utmost earnestness and effort. We request your
approval of this internship report for partial fulfilment of the requirements for the
degree of Bachelor of Science in Computer Science and Engineering.\\[0.3cm]

Thank you.\\[1.5cm]

\begin{minipage}[t]{0.45\textwidth}
Sincerely yours,\\[2cm]
$\overline{\text{Md. Miftahul Islam}}$\\
ID: 221014030\\
Program: BSc.\ CSE\\[1.5cm]
$\overline{\text{Md. Nazirul Islam}}$\\
ID: 221014029\\
Program: BSc.\ CSE
\end{minipage}%
\hfill%
\begin{minipage}[t]{0.45\textwidth}
Academic Supervisor:\\[2cm]
$\overline{\text{Dr.\ Muhammad Golam Kibria, SMIEEE}}$\\
Professor \& Head (Academic)\\
Department of CSE\\
University of Liberal Arts Bangladesh
\end{minipage}

\restoregeometry
\cleardoublepage

% ── Certificate ────────────────────────────────────────────────────────────────
\thispagestyle{empty}

\section*{Certificate}

This is to certify that the report titled \textit{Frontend Architecture and System Design
Program} has been submitted by \textbf{Md. Miftahul Islam} (ID: 221014030) and
\textbf{Md. Nazirul Islam} (ID: 221014029) in partial fulfillment of the requirements
for the Bachelor degree (BSc.\ in Computer Science and Engineering) in the Department
of Computer Science and Engineering of University of Liberal Arts Bangladesh. This is a
record of work carried out by the students under our supervision. No part of the
internship report has been submitted for any other degree and the performance of the
students has been accepted as satisfactory.\\[1.5cm]

\textbf{Supervisor (Academic):}\\[2cm]
$\overline{\text{Dr.\ Muhammad Golam Kibria, SMIEEE}}$\\
Professor \& Head (Academic)\\
Department of Computer Science and Engineering\\
University of Liberal Arts Bangladesh\\[1.5cm]

\textbf{Supervisor (Industry):}\\[2cm]
$\overline{\text{Md.\ Rafshanul Hoque Siam}}$\\
Lead Software Engineer\\
conneqtedagents.ai

\thispagestyle{empty}
\cleardoublepage

% ── Front Matter ───────────────────────────────────────────────────────────────
\pagestyle{phdthesis}
\frontmatter
\begin{doublespace}

\addcontentsline{toc}{chapter}{Executive Summary}
\chapter*{Executive Summary}

This report presents a record of the work we completed during our internship at
\textbf{conneqtedagents.ai} under the supervision of \textbf{Md.\ Rafshanul Hoque
Siam}, Lead Software Engineer. Rather than treating the internship only as an academic
requirement, we approached it as an opportunity to understand how real client work is
planned, reviewed, and delivered inside a professional frontend engineering team.

Our assigned work centered on \textbf{Frontend Architecture and System Design} for a
premium leather manufacturing company's website. The project involved building a modern,
responsive, multi-page web experience using \textbf{Next.js 15 (App Router)},
\textbf{TypeScript}, \textbf{React 19}, \textbf{Tailwind CSS 4}, \textbf{Framer Motion},
and \textbf{Swiper}. We worked on reusable section components, animated page elements,
layout composition, responsive behavior, and content presentation aligned with the
client's visual identity.

This report explains the company context, the purpose of the internship, the methods we
followed during development, and the product outcomes we contributed to. It also records
the technical and professional lessons we gained from the experience and concludes with
recommendations based on the challenges and observations we made during the internship.

\cleardoublepage

\addcontentsline{toc}{chapter}{Acknowledgements}
\chapter*{Acknowledgements}

We would like to express our sincere gratitude to our academic supervisor,
\textbf{Dr.\ Muhammad Golam Kibria, SMIEEE}, Professor \& Head (Academic) of the
Department of Computer Science and Engineering, University of Liberal Arts Bangladesh,
for his invaluable guidance, continuous encouragement, and support throughout the
duration of this internship and in the preparation of this report.

We would also like to extend our heartfelt thanks to our industry supervisor,
\textbf{Md.\ Rafshanul Hoque Siam}, Lead Software Engineer at conneqtedagents.ai,
for his tremendous support, technical mentorship, and hands-on guidance during our
internship. His expertise in frontend architecture and system design greatly enriched
our learning experience.

Finally, we would like to thank everyone at conneqtedagents.ai who provided feedback,
shared knowledge, and made our internship a truly rewarding and growth-oriented experience.

\cleardoublepage

\begin{spacing}{1.5}
\tableofcontents
\end{spacing}

\cleardoublepage

\addcontentsline{toc}{chapter}{List of Figures}
\listoffigures

\cleardoublepage

% ── Main Matter ────────────────────────────────────────────────────────────────
\mainmatter

\chapter{Introduction}

\section{Presentation of the Company}

\begin{flushleft}
The internship program gave us direct exposure to the professional environment that exists
beyond classroom projects. Through this placement, we were able to observe how actual
client requirements are translated into engineering tasks, reviewed through collaboration,
and refined into production-ready frontend deliverables.

conneqtedagents.ai is a software company that works on modern web products and intelligent
digital systems for local and international clients. From our experience inside the team,
the company places strong emphasis on structured frontend architecture, maintainable code,
practical design execution, and the steady use of modern development tools.

What stood out most during our internship was the company's working style: tasks were
handled through iterative delivery, review-based improvement, and close attention to both
technical correctness and user-facing quality. This environment gave us a realistic view
of how a professional software team balances speed, clarity, and product standards.
\end{flushleft}

\section{Objective of the Internship}

The objective of this internship was not only to satisfy the academic requirement of our
Bachelor degree, but also to gain experience in the kind of engineering work that cannot
be fully reproduced in coursework alone. We wanted to understand how frontend systems are
actually designed, reviewed, and delivered when the output is meant for a real client and
real users.

By joining conneqtedagents.ai as frontend engineering interns, we aimed to strengthen our
skills in component-based development, responsive design, animation-heavy user interfaces,
and structured application architecture. We were also interested in learning how senior
engineers make decisions about code quality, maintainability, visual polish, and delivery
workflow inside a production setting.

For us, the internship served as a transition point between academic learning and
professional practice. It allowed us to test our existing knowledge, identify our weak
areas, and build confidence through meaningful project work rather than isolated exercises.

\section{Scope and Methodology}

\subsection{Overview}

This report provides a comprehensive account of our internship experience at
conneqtedagents.ai during the Spring 2026 semester. The scope of the internship covered
the full frontend development lifecycle of a premium leather company portfolio website,
from initial requirement analysis and component architecture design through to
implementation, testing, and final delivery.

Our work was scoped to the frontend layer of the product --- encompassing page layout,
section component development, animation implementation, responsive design, and
client-side interactivity. The methodology we followed was aligned with professional
software engineering practices used in production environments, adapted to the specific
demands of a visually driven, multi-page marketing website.

\subsection{Technology Stack}

The project was built on a modern, production-ready frontend technology stack selected
by the organization to ensure performance, scalability, and developer experience:

\begin{itemize}
  \item \textbf{Next.js 15 (App Router)} for file-based routing, server and client
  component management, metadata handling, and static page generation.
  \item \textbf{TypeScript} for static type safety, precise interface definitions, and
  compile-time error detection across the entire codebase.
  \item \textbf{React 19} as the core UI rendering library, used with functional
  components and hooks throughout.
  \item \textbf{Tailwind CSS 4} as the utility-first styling framework for building
  a fully responsive and visually consistent design system.
  \item \textbf{Framer Motion} for scroll-driven parallax effects, entrance animations,
  staggered transitions, and page-level shutter animations.
  \item \textbf{Swiper} for touch-friendly, responsive image carousel and product
  gallery components.
  \item \textbf{DOMPurify \& jsdom} for sanitizing rich HTML content rendered in
  text-heavy sections, preventing XSS vulnerabilities.
  \item \textbf{Lucide React \& React Icons} for scalable, consistent iconography
  across navigation, call-to-action, and informational components.
\end{itemize}

\subsection{Development Process}

Our development process followed an iterative, task-driven workflow aligned with the
professional conventions of the organization:

\begin{itemize}
  \item \textbf{Requirement Understanding:} Each task began with a review of design
  specifications and a discussion with our industry supervisor to clarify expected
  behavior, visual standards, and technical constraints.
  \item \textbf{Branch-Based Development:} We worked on dedicated feature branches for
  each section or component, keeping changes isolated and enabling clean pull request
  reviews.
  \item \textbf{Implementation:} Features were implemented following the project's
  established folder structure and naming conventions, ensuring consistency with the
  existing codebase.
  \item \textbf{Code Review and Iteration:} Every pull request was reviewed by our
  industry supervisor before merging. Review feedback was addressed promptly, and
  revisions were made in the same branch before re-submission.
  \item \textbf{Delivery and Integration:} Completed and approved features were merged
  into the main branch and verified in the integrated build before being considered
  delivered.
\end{itemize}

\subsection{Component Design and Responsiveness}

A core principle of our development approach was component-driven architecture.
Each website section was encapsulated as a self-contained React component with
well-defined, typed props. This enabled reuse across multiple pages and made
individual sections easy to modify or extend without side effects.

Responsiveness was treated as a first-class requirement. We used Tailwind CSS
responsive modifiers (\texttt{sm:}, \texttt{md:}, \texttt{lg:}) to build layouts
that adapt fluidly across mobile, tablet, and desktop breakpoints. Key patterns
included collapsible navigation for mobile, flexible two-column content-image
sections that stack on smaller screens, and grid-based category layouts with
dynamic column counts.

Page-level composition was handled by page components under \texttt{components/pages/},
which assembled the correct section components for each route. This separation between
layout composition and section implementation kept each file focused and maintainable.

\subsection{Animation and Integration}

Animations were implemented using Framer Motion throughout the project. Scroll-driven
effects were achieved via the \texttt{useScroll} and \texttt{useTransform} hooks, which
allowed background parallax and content fade effects to respond continuously to the
user's scroll position. Entrance animations for page sections used \texttt{motion}
variants with \texttt{initial}, \texttt{animate}, and \texttt{exit} states, while the
\texttt{AnimatePresence} component managed the lifecycle of conditionally rendered
elements such as the mobile navigation drawer.

The page-level shutter transition was implemented as a higher-order component
(\texttt{withShutter}), wrapping each page with a branded entrance animation on
initial load. This approach kept animation logic fully decoupled from page content
components.

Rich HTML content, sourced from mock data and intended for eventual CMS integration,
was rendered using \texttt{dangerouslySetInnerHTML} with DOMPurify sanitization applied
before injection, ensuring that content-driven sections remained safe against
cross-site scripting vulnerabilities.

\subsection{Testing, Documentation, and Reporting}

Given the nature of the project as a frontend portfolio website, testing focused on
functional and visual correctness rather than automated unit tests. Each completed
component was manually tested across multiple viewport sizes using the browser's
responsive design mode, verifying layout integrity at mobile, tablet, and desktop
breakpoints.

Interactive components --- including the navbar drawer, shutter transition, parallax
hero, and image stack previewer --- were tested for correct animation behavior, smooth
transitions, and the absence of visual artifacts or layout shifts.

Code documentation was maintained through clear component prop interfaces using
TypeScript, which served as inline documentation for expected inputs, types, and
optional configurations. Shared types were centralized in the \texttt{interfaces/}
directory, making them discoverable and reusable across the codebase.

Progress reporting was conducted through regular check-ins with our industry supervisor,
during which completed sections were reviewed visually, feedback was captured, and
upcoming tasks were prioritized. This internship report itself constitutes the formal
written record of the work carried out, the methodologies applied, and the outcomes
achieved during the Spring 2026 internship period.

\chapter{conneqtedagents.ai}

\section{About conneqtedagents.ai}

Based on our experience during the internship, conneqtedagents.ai can be described as a
modern software company that builds digital products with a strong focus on structured
engineering practices and real business usefulness. The company works across several kinds
of web-based solutions, including client-facing websites, internal systems, and intelligent
software products that incorporate contemporary frontend and backend technologies.

From the engineering culture we observed, the organization values maintainability,
incremental improvement, and clarity in implementation. The work environment encouraged us
to think beyond simply making interfaces look correct; we were expected to consider how
components scale, how design decisions affect future maintenance, and how code should be
organized so that it remains understandable to other developers.

\section{Mission and Vision}

\noindent\underline{Mission}: From our understanding of the company's work, its mission is
to deliver practical and well-engineered digital products that help businesses solve
real operational and communication problems through technology.

\noindent\underline{Vision}: Its long-term direction appears to be centered on building
reliable, modern, and increasingly intelligent software systems that combine strong user
experience with scalable technical foundations.

\section{Core Values and Engineering Culture}

conneqtedagents.ai maintains a strong engineering culture centered on:

\begin{itemize}
  \item \textbf{Quality:} Every product is built with a commitment to code quality,
  performance, and reliability.
  \item \textbf{Innovation:} The team constantly adopts and evaluates emerging
  technologies to stay ahead of industry trends.
  \item \textbf{Collaboration:} Open communication and iterative feedback are core
  to how the team operates.
  \item \textbf{Developer Experience:} The company invests in tooling, workflows, and
  processes that empower engineers to do their best work.
\end{itemize}

\section{Services}

\noindent\textbf{i.\ AI-Powered Platform Development}

The company develops intelligent platforms that combine modern web technologies with
automation-oriented features to support business workflows and digital operations.

\noindent\textbf{ii.\ Frontend Architecture and System Design}

conneqtedagents.ai also provides frontend engineering services focused on reusable
component architecture, performance-conscious implementation, and polished interfaces
built with tools such as Next.js, TypeScript, and React.

\noindent\textbf{iii.\ Full-Stack Web Application Development}

The company works across full-stack product development as well, covering both backend
service design and user-facing application layers.

\noindent\textbf{iv.\ Cloud-Native Deployment and DevOps}

It also follows modern deployment and delivery practices that support maintainable
software releases and long-term product operation.

\chapter{Details of Internship Work}

The main objective of the internship was to gain experience in a real professional
software development environment and to apply the knowledge acquired during our academic
program to production-grade engineering tasks. Throughout our internship at
conneqtedagents.ai, we were actively involved in the design and development of a
\textbf{premium leather company portfolio website} --- a fully responsive, multi-page
marketing website built using \textbf{Next.js 15 (App Router)}, \textbf{TypeScript},
\textbf{React 19}, \textbf{Tailwind CSS 4}, \textbf{Framer Motion}, and \textbf{Swiper}.

The website serves as the digital presence of a premium leather manufacturing brand,
showcasing the company's product catalogue, leather craftsmanship heritage, service
offerings, and sustainability commitments. The project required careful attention to
visual quality, animation fidelity, and component reusability across multiple pages.

\section{Technologies Used}

\begin{itemize}
  \item \textbf{Next.js 15 (App Router):} A React-based framework used for file-based
  routing, server and client component separation, metadata management, and static
  page generation. The App Router paradigm enabled clean layout composition across
  shared navigation and footer structures.
  \item \textbf{TypeScript:} A strongly typed superset of JavaScript that enabled us
  to write safer, more maintainable code with compile-time type checking, precise
  interface definitions, and generic component props across the entire codebase.
  \item \textbf{React 19:} The core UI library used for building component-based user
  interfaces. We used React hooks for managing component state, scroll effects, and
  animated transitions.
  \item \textbf{Tailwind CSS 4:} A utility-first CSS framework used to design a
  fully responsive and visually consistent UI. We leveraged Tailwind's composition
  model to build reusable, maintainable styles across all page sections.
  \item \textbf{Framer Motion:} A production-ready animation library for React, used
  extensively to implement parallax hero effects, scroll-driven opacity transitions,
  staggered list animations, and page-level shutter transitions.
  \item \textbf{Swiper:} A modern touch-friendly slider library used for implementing
  product and image gallery carousels with smooth swipe gestures and responsive
  breakpoints.
  \item \textbf{Lucide React \& React Icons:} Icon libraries used to provide
  consistent, scalable iconography across navigation menus, call-to-action buttons,
  and informational sections.
  \item \textbf{DOMPurify \& jsdom:} Used to safely sanitize server-rendered HTML
  strings injected into rich-text section components, preventing XSS vulnerabilities
  in content-driven sections.
\end{itemize}

\section{Project Architecture}

The project followed a structured folder architecture aligned with Next.js App Router
conventions:

\begin{itemize}
  \item \textbf{app/:} Contains all route segments, each with a dedicated
  \texttt{page.tsx} file. Routes include \texttt{/}, \texttt{/about-us},
  \texttt{/products}, \texttt{/services}, \texttt{/leather-experience},
  \texttt{/contact-us}, \texttt{/privacy-policy}, and \texttt{/cookie-policy}.
  \item \textbf{components/sections/:} Reusable section-level components such as the
  hero section, category layout, content-image section, footer, and navbar.
  \item \textbf{components/pages/:} Page-level composition components that assemble
  the appropriate section components for each route.
  \item \textbf{components/UI/:} Atomic UI primitives shared across sections.
  \item \textbf{assets/:} Centralized image asset exports organized by category
  (products, categories, logo, others).
  \item \textbf{constants/, mock/, interfaces/, utils/:} Shared data, TypeScript
  type definitions, mock content, and utility functions.
\end{itemize}

\section{Key Features Developed}

During our internship, we contributed to the design and implementation of all major
website sections. The following subsections describe the key features delivered.

\subsection{Animated Hero Sections with Parallax Effect}

We implemented a full-screen hero section with a parallax background effect using
Framer Motion's \texttt{useScroll} and \texttt{useTransform} hooks. The background
image translates vertically as the user scrolls, while the foreground content fades
out smoothly. This section is reused across all major pages with configurable title,
description, and optional breadcrumb navigation.

\begin{figure}[htbp]
  \centering
  \includegraphics[width=15cm, height=9cm]{hero_section.png}
  \caption{Animated Hero Section with Parallax Scroll Effect}
  \label{fig:hero_section}
\end{figure}

\newpage

\subsection{Category Layout and Product Showcase}

We built a category layout section that displays the leather product categories
(Nautical, Aviation, Fashion, Automotive, Interior Design, and Regenerated Materials)
in a visually rich grid. Each category card is backed by high-resolution \texttt{.avif}
images and includes a title, description, and navigation link. The section integrates
a description header built from a shared \texttt{DescriptionSection} component.

\begin{figure}[htbp]
  \centering
  \includegraphics[width=15cm, height=9cm]{category_layout.png}
  \caption{Leather Category Layout Section}
  \label{fig:category_layout}
\end{figure}

\newpage

\subsection{Content-Image Sections}

We implemented a flexible \texttt{ContentImageSection} component that renders a
two-column layout alternating between text content and an image. The layout direction
is configurable (\texttt{leftToRight} or \texttt{rightToLeft}), supporting reuse across
multiple pages with different content and visual orientations. Rich HTML content is
safely rendered using DOMPurify sanitization.

\begin{figure}[htbp]
  \centering
  \includegraphics[width=15cm, height=9cm]{content_image_section.png}
  \caption{Configurable Content-Image Section}
  \label{fig:content_image}
\end{figure}

\newpage

\subsection{Lean Image Stack Previewer}

We developed a lean image stack previewer component that displays a layered stack of
product images with a hover-driven reveal interaction. This component was used on the
home page to provide an interactive visual preview of the company's leather product
range, enhancing engagement without requiring a full carousel.

\begin{figure}[htbp]
  \centering
  \includegraphics[width=15cm, height=9cm]{image_stack_previewer.png}
  \caption{Lean Image Stack Previewer Component}
  \label{fig:image_stack}
\end{figure}

\newpage

\subsection{Animated Navbar with Drawer Navigation}

We implemented a responsive navbar component with scroll-aware styling. On scroll, the
navbar transitions from a transparent overlay to a solid background using Framer Motion
animation controls. On mobile, a full-screen animated drawer menu slides in with
staggered link animations, providing a smooth and accessible navigation experience.
The drawer includes page descriptions alongside each navigation link.

\begin{figure}[htbp]
  \centering
  \includegraphics[width=15cm, height=9cm]{navbar_drawer.png}
  \caption{Responsive Navbar with Animated Drawer}
  \label{fig:navbar}
\end{figure}

\newpage

\subsection{Page Transition Shutter HOC}

We implemented a higher-order component (HOC) called \texttt{withShutter} that wraps
page components with an animated entrance transition. On initial render, a shutter-style
overlay sweeps across the screen and reveals the page content, providing a polished
brand-consistent page load experience across all routes.

\begin{figure}[htbp]
  \centering
  \includegraphics[width=15cm, height=9cm]{shutter_transition.png}
  \caption{Shutter Page Transition Higher-Order Component}
  \label{fig:shutter}
\end{figure}

\chapter{Key Learnings}

From both a technical and professional perspective, the internship at conneqtedagents.ai
proved to be an immensely valuable learning experience. Working on real client-facing
frontend products gave us practical insight into how projects are analyzed, planned,
implemented, reviewed, and delivered in a professional environment.

\section{Analysis and Findings}

During the internship, we observed that successful frontend delivery depends not only on
writing functional code, but also on understanding user needs, aligning with branding
requirements, and maintaining consistency across the entire product surface. Our analysis
of the assigned work showed that reusable component architecture, responsive design, and
carefully managed animation patterns were essential for sustaining development speed and
visual quality.

We also found that a project with multiple content-heavy pages benefits significantly
from centralized assets, shared section components, and typed interfaces. These patterns
reduce duplication, simplify maintenance, and improve collaboration between developers.

\section{Project Outcomes}

The internship resulted in the successful delivery of multiple frontend outputs that
demonstrated both technical execution and practical product value.

\subsection{Personal Portfolio / Company Website}

One of the major outcomes of the internship was the design and development of a
multi-page company website for a premium leather manufacturing brand. The website
showcases the company's services, products, brand identity, and leather craftsmanship
through modern layouts, animated hero sections, content-image compositions, and a
responsive navigation system.

This outcome demonstrated our ability to build a polished public-facing product using
modern frontend technologies while maintaining reusability, responsiveness, and visual
consistency across pages.

\subsection{Custom Admin Panel (ELF Dashboard)}

In addition to the company website work, the internship also exposed us to the broader
expectations of internal product interfaces such as custom dashboards and structured
admin-facing systems. Through this context, we developed a clearer understanding of how
frontend architecture differs between marketing websites and dashboard-oriented products.

We learned that an admin panel such as an ELF Dashboard requires a stronger emphasis on
information hierarchy, reusable data display patterns, role-aware navigation, and task
efficiency. Even where our direct contribution was centered on the website, the exposure
to this category of product improved our understanding of scalable frontend system design.

\section{Learning and Development}

The internship accelerated both our technical growth and our professional maturity.
The following subsections summarize the main areas of learning.

\subsection{Technical Learning}

The technical learning gained during the internship was substantial. We acquired
hands-on experience with:

\begin{itemize}
  \item \textbf{Next.js 15 App Router:} We learned how to structure route segments,
  separate server and client components, manage layouts, and configure page-level
  metadata in a production codebase.
  \item \textbf{TypeScript in Production:} We practiced writing interfaces and typed
  props for reusable components, improving safety, readability, and maintainability.
  \item \textbf{Component-Driven Architecture:} We learned to build section-level and
  page-level components that can be reused across different routes without unnecessary
  duplication.
  \item \textbf{Responsive Design:} We deepened our ability to build interfaces that
  adapt cleanly across mobile, tablet, and desktop devices using Tailwind CSS.
  \item \textbf{Animation Systems:} We became proficient with Framer Motion for
  parallax effects, transitions, and motion-driven interactions that improve the user
  experience.
  \item \textbf{Safe Content Rendering:} We learned the importance of sanitizing HTML
  content using DOMPurify before rendering it in the DOM, helping prevent XSS-related
  vulnerabilities.
\end{itemize}

\subsection{Professional Learning}

The internship also contributed significantly to our professional development.

\begin{itemize}
  \item We learned how to communicate progress clearly during supervisor check-ins.
  \item We developed discipline in following project conventions and maintaining code
  consistency within an existing codebase.
  \item We learned to accept and apply code review feedback constructively.
  \item We improved our ability to estimate work, break down tasks, and iterate on
  features based on feedback.
  \item We gained a practical understanding of how client expectations influence design,
  content presentation, and delivery quality.
\end{itemize}

\section{Tools and Technologies Used}

The internship required consistent use of a modern frontend toolchain. For clarity,
the technologies can be grouped as follows.

\subsection{Frontend}

The frontend layer was primarily built using \textbf{Next.js 15}, \textbf{React 19},
\textbf{TypeScript}, and \textbf{Tailwind CSS 4}. These technologies formed the
foundation for routing, rendering, component design, styling, and application structure.

\subsection{Libraries and Tools}

To enhance functionality and developer productivity, we worked with several additional
libraries and tools including \textbf{Framer Motion} for animation, \textbf{Swiper}
for image sliders, \textbf{Lucide React} and \textbf{React Icons} for iconography,
and \textbf{DOMPurify} with \textbf{jsdom} for safe HTML sanitization.

\subsection{Development Environment and Version Control}

The development environment centered around \textbf{Visual Studio Code}, where we used
extensions and integrated terminal workflows for coding, previewing, and debugging the
application. Source code was managed using \textbf{Git}, with branch-based development
and pull request workflows supporting collaboration, review, and controlled integration
of completed features.

\section{User Guidelines}

\subsection{Company Website}

The company website is intended to present the brand, products, and services in a clear
and visually engaging manner. Users should be able to navigate between the Home,
About Us, Products, Services, Leather Experience, Contact Us, Privacy Policy, and
Cookie Policy pages using the primary navigation menu. On mobile devices, the drawer
navigation provides the same access in a compact form.

Users can browse category-focused sections, read descriptive content about the
company's philosophy and craftsmanship, and interact with animated content blocks
designed to improve engagement without reducing clarity.

\subsection{Admin Dashboard}

For admin-oriented products such as a custom dashboard, the key guideline is to design
interfaces around clarity, efficiency, and task completion. Users should be able to
identify navigation options quickly, locate data without confusion, and perform common
actions with minimal friction. From our internship experience, we learned that admin
interfaces must prioritize usability, consistency, and structured information layout
over purely decorative presentation.

\section{Summary of Findings}

In summary, our internship showed us that frontend engineering is not just about writing
components that render correctly on screen. It involves making decisions about structure,
clarity, responsiveness, animation, maintainability, and user expectations at the same
time. The strongest learning came from handling real project constraints, where every
technical choice had visible impact on delivery quality.

The experience made us more aware of how professional frontend work is evaluated in
practice. A successful result depends on technical accuracy, visual consistency,
communication with reviewers, and the ability to revise work without losing overall
system coherence. This understanding has prepared us more realistically for future roles
in frontend and product-focused software development.

\chapter{Conclusion}

\section{Major Benefits of the Internship}

The internship at conneqtedagents.ai was an exceptionally rewarding experience that
delivered significant benefits both technically and professionally. Working on a
production-grade portfolio website for a premium leather manufacturing company gave us
direct exposure to the standards and expectations of real client delivery in a
professional software environment.

From a technical standpoint, we gained hands-on proficiency in Next.js 15 (App Router),
TypeScript, React 19, Tailwind CSS 4, Framer Motion, and Swiper --- a modern and
commercially relevant technology stack. We moved beyond academic exercises to building
real animated components, responsive layouts, and reusable architectural patterns that
are deployed for an actual client.

From a professional standpoint, the internship taught us how to interpret client
requirements, collaborate under the guidance of a senior engineer, manage iterative
delivery, and participate in code review cycles. We also developed a stronger sense of
responsibility for code quality, security, and performance in a real product context.

\section{Limitations of the Internship}

The internship period, while highly productive, was not without its challenges and
limitations. The primary challenge we encountered was adapting to the design language
and visual standards expected by an international client. The website demanded a high
level of visual polish, which required a more design-oriented mindset than typical
application development exercises in an academic setting.

Several tools and techniques --- including advanced Framer Motion scroll-driven
animations, Swiper's responsive breakpoint configuration, and DOMPurify-based HTML
sanitization --- were new to us at the start and required independent research before
we could apply them confidently.

Performance considerations for image-heavy pages, using \texttt{.avif} assets and
external image sources, also required careful attention. Time constraints limited the
extent to which we could optimize image loading strategies, apply lazy loading
throughout, and conduct formal performance auditing using Lighthouse.

Furthermore, the internship was primarily frontend-focused, which limited our exposure
to backend API design, server deployment workflows, and continuous integration pipelines
--- areas that are equally important in full-stack product development.

\section{Recommendations for conneqtedagents.ai}

Based on our experience during the internship, we offer the following recommendations
to conneqtedagents.ai for improving future intern engagements and project quality:

\begin{itemize}
  \item \textbf{Structured Onboarding:} Preparing a written onboarding guide covering
  the project's design system, component conventions, animation patterns, and folder
  architecture would significantly reduce the initial ramp-up time for new contributors.
  \item \textbf{Design System Documentation:} Formalizing the typography scale, spacing
  conventions, color palette, and animation timing guidelines into a documented design
  system would help maintain visual consistency as the project scales.
  \item \textbf{Image Optimization Pipeline:} Switching to Next.js \texttt{<Image />}
  components throughout and serving all assets from a CDN would meaningfully improve
  Lighthouse performance scores and Core Web Vitals.
  \item \textbf{Component Library:} Introducing a component storybook or visual
  regression testing suite would support long-term maintainability and prevent visual
  regressions as new sections are added.
\end{itemize}

\section{Recommendations for Future Interns}

We offer the following recommendations to students who will undertake a similar
internship in frontend engineering:

\begin{itemize}
  \item \textbf{Study the App Router paradigm early:} The distinction between server
  and client components in Next.js 15 is fundamental. Understanding this model before
  starting will allow you to contribute more confidently from the first week.
  \item \textbf{Practice animation-driven development:} Framer Motion is a powerful
  but nuanced library. Spending time with its documentation and sandbox examples prior
  to the internship will reduce friction when implementing complex transitions.
  \item \textbf{Understand security basics:} Be aware of XSS vulnerabilities when
  rendering dynamic HTML. Understanding how and why to use DOMPurify is a practical
  skill that will be relevant on almost any content-driven project.
  \item \textbf{Embrace code review:} Treat code review not as a quality gate but as
  a learning channel. The feedback received during reviews is often the fastest path
  to improving your engineering judgment.
  \item \textbf{Ask early, iterate often:} Clarifying design and requirement ambiguities
  with your supervisor at the start of each task prevents costly late-stage revisions
  and builds trust with the team.
\end{itemize}

\section{Recommendations for the University and Internship Program}

Based on our internship experience, we also offer recommendations directed at the
University of Liberal Arts Bangladesh and the CSE internship program:

\begin{itemize}
  \item \textbf{Introduce modern frontend tooling in the curriculum:} Courses
  covering Next.js, TypeScript, and component-driven architecture would better prepare
  students for the kinds of tasks they encounter in modern frontend internships.
  \item \textbf{Include animation and UI libraries in elective offerings:} Tools such
  as Framer Motion and Tailwind CSS are industry-standard but absent from most
  traditional curricula. Exposure to these early would reduce the learning curve
  during internships.
  \item \textbf{Strengthen the industry-academia bridge:} More frequent interaction
  between industry supervisors and the university throughout the internship period
  would help align expectations and ensure that students are being exposed to
  educationally valuable tasks.
  \item \textbf{Encourage portfolio-based internship placements:} Internships that
  result in a publicly accessible deliverable (such as a deployed website) provide
  students with tangible evidence of professional work, which is highly valuable for
  early career placement.
\end{itemize}

\section{Final Remark}

The internship at conneqtedagents.ai became one of the most meaningful parts of our
undergraduate journey because it required us to apply what we had learned in a context
where quality, timing, review feedback, and client expectations all mattered at once.
That shift from academic exercise to professional responsibility changed how we now think
about software work.

We are ending this internship with stronger technical judgment, better work discipline,
and a much clearer understanding of the level of care required in modern frontend
development. We remain sincerely grateful to our academic supervisor
\textbf{Dr.\ Muhammad Golam Kibria, SMIEEE}, our industry supervisor
\textbf{Md.\ Rafshanul Hoque Siam}, and the team at conneqtedagents.ai for trusting us
with meaningful work and helping us grow through that process.

\end{doublespace}

\backmatter
\cleardoublepage

\addcontentsline{toc}{chapter}{References}
\chapter*{References}

\begin{itemize}
  \item Next.js Documentation\\
  \url{https://nextjs.org/docs}

  \item TypeScript Handbook\\
  \url{https://www.typescriptlang.org/docs/handbook/}

  \item React Documentation\\
  \url{https://react.dev/}

  \item Tailwind CSS Documentation\\
  \url{https://tailwindcss.com/docs}

  \item Framer Motion Documentation\\
  \url{https://www.framer.com/motion/}

  \item Swiper Documentation\\
  \url{https://swiperjs.com/}

  \item DOMPurify --- XSS Sanitizer\\
  \url{https://github.com/cure53/DOMPurify}

  \item Lucide React Icons\\
  \url{https://lucide.dev/}

  \item conneqtedagents.ai\\
  \url{https://conneqtedagents.ai}
\end{itemize}

\bibliographystyle{IEEEtran}
\cleardoublepage
\addcontentsline{toc}{chapter}{Bibliography}
\bibliography{bibs}

\end{document}
